\section{Lecture Principles}

\subsection{Multi-level Testing in AI Research}

% Multi-level testing is common in software engineering but less systematic in traditional AI research. In software engineering, testing progresses from individual functions to integrated components, the complete system, and finally user-level validation. This structure is important for production software because the final system must operate reliably for customers. In contrast, many AI research projects are developed mainly for experimentation and are used only by researchers. Therefore, unusual edge cases may be less important than ensuring that the main data-processing, training, and evaluation pipelines correctly implement the research idea.
Multi-level testing is common in software engineering, but it is less systematic in traditional AI research. In software engineering, testing typically starts with individual functions, gradually expanding to integrated components, the entire system, and finally user-level verification. This structure is crucial for production software because the final system must be able to run reliably for users. In contrast, many AI research projects are primarily experimental, intended only for researchers. To turn the AI research into a product would be another project. Therefore, unusual edge cases may be less important than ensuring that the main data-processing, training, and evaluation pipelines correctly implement the research idea.

However, I have found testing increasingly valuable as AI coding agents, such as Codex and Claude Code, become more involved in my project. These agents can generate substantial parts of functions and pipelines, but the generated code may not always reflect my intended behaviour. Tests provide a practical way to control its quality. For example, unit tests can check individual preprocessing or graph-construction functions, while integration tests can confirm that the data pipeline and model work together correctly. In this setting, tests also act as executable specifications of my ideas. They help verify that the agent has implemented what I intended. AI agents can reduce implementation effort, but they do not remove the need to judge whether generated code is correct. Therefore, multi-level testing is a useful software-engineering principle that can be borrowed to improve the reliability of AI research code in the era of AI coding agents.


\subsection{The Science-Engineering Loop in my Research}

My research project combines scientific investigation with substantial engineering work. From a scientific perspective, I investigate which methods can best learn representations of football player motion from tracking data. However, unlike research based on standardized and preprocessed benchmarks, my work starts from raw data.
As a result, a lot of engineering work is required, including processing the raw tracking data, filtering and grouping samples, constructing datasets, creating annotations, and defining appropriate data splits. To support this process, I also developed an annotation software system. These parts must be reliable enough to ensure that the experimental results are trustworthy.

The lecture describes machine learning as both an empirical science and a design discipline, involving an iterative cycle of experimentation, implementation, observation, and refinement. This closely reflects my doctoral project. Although the main goal is to produce strong scientific contributions, the validity of those contributions depends heavily on the quality of the underlying engineering. Errors in data processing, annotation, dataset construction, or evaluation could lead to incorrect scientific conclusions. Therefore, my research is scientific in its objectives but strongly dependent on engineering in its execution.


